% ===========================================================================
\chapter{Considerações Finais}\label{cap:consideracoes}
% ===========================================================================

Este capítulo faz o balanço do trabalho em relação aos objetivos do
\autoref{cap:introducao}, aponta o que é contribuição dos autores e indica
caminhos para continuar o projeto.

\section{Conclusões do Projeto de Formatura}

O objetivo geral foi atingido. O protótipo vende livros digitais organizados em
grupos de mesmo preço, entrega ao cliente as chaves dos N livros escolhidos por
meio de \emph{Oblivious Transfer} e distribui os livros cifrados por uma CDN,
de modo que a loja recebe o pagamento sem saber quais livros foram comprados.
O fluxo roda de ponta a ponta com os três componentes previstos, e todos os
requisitos funcionais essenciais foram implementados e verificados pela
bateria controlada da \autoref{sec:testes}.

Quanto aos objetivos específicos:

\begin{itemize}
\item \textbf{Arquitetura:} a separação entre sistema web, CDN e programa local
foi mantida do início ao fim. A principal revisão foi permitir que o programa
local acesse a rede, restrito ao sistema web e à CDN, o que eliminou o
transporte manual de arquivos sem enfraquecer a garantia sobre a escolha,
porque o código do receptor continua fora do controle do sistema web.

\item \textbf{Algoritmo criptográfico:} o OT 1-out-of-n de
\citeonline{naorpinkas1999}, repetido N vezes sobre o OT base de
\citeonline{masnyrindal2019}, foi integrado por meio da biblioteca libOTe. O
custo, de N vezes $\lceil \log_2 n \rceil$ OTs base por compra, mostrou-se
pequeno o bastante para dispensar OT \emph{Extension} nos tamanhos de grupo do
protótipo.

\item \textbf{Legitimidade sem revelar o conteúdo:} a cobrança pelo grupo e
pela quantidade, os tokens de uso único com assinatura cega e o reporte da CDN
apenas com tokens e quantidade permitem cobrar o cliente e remunerar a CDN sem
que o sistema web aprenda os livros comprados.

\item \textbf{Testes de segurança e desempenho:} os testes mostram que o que o
sistema web grava, registra e responde não varia com a escolha, e que nenhuma
chave sai do computador do cliente. As metas de tempo de busca e de compra
foram atendidas com folga, e a de vazão foi atendida a partir de quatro
núcleos.
\end{itemize}

Alguns resultados não foram atingidos ou ficaram parciais. A disponibilidade
(RNF06), a taxa de compras concluídas sem suporte (RNF07) e o tempo de
manutenção (RNF08) dependem de operação real e de estudo com usuários, que
ficaram fora do alcance de um protótipo. O pagamento usa um \emph{gateway}
simulado, porque um real exige conta de lojista, e a janela gráfica existe
apenas para Windows. A defesa contra o ataque de falha seletiva no OT é uma
escolha de engenharia, e não uma garantia formal. Por fim, a não conivência
entre o sistema web e a CDN continua sendo uma premissa: o sistema reduz o que
cada parte aprende, mas não impede que as duas juntem o que sabem.

\section{Contribuições}

A literatura sobre OT com preço e transações oblívias concentra-se nos
protocolos e em suas provas. A contribuição deste trabalho é de engenharia: um
sistema completo, que integra o OT a pagamento, entrega e remuneração, e a
análise do que cada decisão preserva ou enfraquece. São de autoria da equipe:

\begin{itemize}
\item a arquitetura em três componentes com separação de conhecimento entre
sistema web, CDN e \emph{gateway}, e a tabela do que cada parte pode saber,
usada para revisar o código e o esquema do banco;
\item a simplificação do OT com preço por grupos de mesmo preço, cobrando pelo
grupo e pela quantidade, com versões imutáveis dos grupos para fixar os
índices do OT;
\item a recuperação de compras interrompidas a partir de uma semente da qual
deriva toda a aleatoriedade do cliente, inclusive o receptor do OT e o
cegamento dos tokens, o que torna útil o envio da resposta por e-mail sem
alterar a biblioteca de OT nem guardar seu estado interno;
\item a remuneração da CDN por tokens de assinatura cega, com uma chave por
faixa de preço publicada no catálogo assinado, e o reporte em lotes
embaralhados, sem livro nem horário;
\item o catálogo como pacote determinístico, assinado e igual para todos, que
leva a navegação e a busca para o computador do cliente e permite detectar um
catálogo personalizado;
\item a implementação dos três componentes, do binário auxiliar sobre o libOTe,
dos testes de privacidade e da bateria controlada de ponta a ponta.
\end{itemize}

As primitivas criptográficas não são de autoria da equipe e vêm de bibliotecas
consolidadas: o OT base do libOTe, a assinatura cega RSA da biblioteca circl e
as cifras, resumos e assinaturas da biblioteca padrão da linguagem Go. A
construção de Naor e Pinkas sobre os OTs base foi implementada pela equipe a
partir do artigo original.

\section{Perspectivas de Continuidade}

\begin{itemize}
\item \textbf{Compromissos das chaves no catálogo:} publicar no pacote
assinado um compromisso da chave de cada livro, para que o programa local
confira as chaves extraídas e a defesa contra a falha seletiva deixe de
depender apenas do comportamento uniforme.

\item \textbf{Conjunto de anonimato e correlação de horários:} estudar o
tamanho mínimo de grupo e de faixa em função do volume de compras, atrasos ou
lotes no download e o uso de redes de anonimato, para reduzir o que se pode
inferir mesmo com conivência entre o sistema web e a CDN.

\item \textbf{Pagamento real com privacidade:} integrar um \emph{gateway} real
e avaliar meios de pagamento que também escondam a identidade do comprador,
tema que ficou fora do escopo.

\item \textbf{Desempenho:} reduzir o custo das assinaturas cegas, que dominam o
tempo da compra, e avaliar OT \emph{Extension} ou OT k-out-of-n para grupos e
quantidades maiores.

\item \textbf{Portabilidade e usabilidade:} levar a janela gráfica a outros
sistemas operacionais, oferecer cópia de segurança cifrada da semente e
conduzir um estudo com usuários para avaliar RNF07.

\item \textbf{Operação:} implantar o sistema web e a CDN em servidores
independentes, operados por partes diferentes, e medir disponibilidade e vazão
em condições reais.
\end{itemize}
